% LSPV-3 REPORT::QGUARD1::2026-10-12T03:36:00Z
\documentclass[11pt]{article}
\usepackage[a4paper,margin=0.9in]{geometry}
\usepackage{amsmath,amssymb,amsthm,booktabs,array,enumitem,xurl,hyperref,microtype}
\hypersetup{colorlinks=true,linkcolor=blue,urlcolor=blue,citecolor=blue}
\setlength{\emergencystretch}{3em}
\newtheorem{theorem}{Theorem}
\newtheorem{corollary}{Corollary}
\newcommand{\EXACT}{\textbf{PROVED/EXACT}}
\newcommand{\KNOWN}{\textbf{KNOWN THEOREM}}
\newcommand{\COMP}{\textbf{COMPUTATIONAL}}
\newcommand{\COND}{\textbf{CONDITIONAL}}
\newcommand{\CAND}{\textbf{CANDIDATE}}
\newcommand{\OPEN}{\textbf{OPEN}}
\newcommand{\RETIRED}{\textbf{FALSIFIED/RETIRED}}
\newcommand{\SAT}{\mathsf{SAT}}
\title{LSPV-3 QGUARD1\\Easy-island avoidance does not negate a worst-case obstruction}
\author{UNIVERSAL\_CORE\_REDTEAM}
\date{12 October 2026 (UTC)}
\begin{document}
\maketitle

\begin{center}
\fbox{\parbox{0.94\linewidth}{
\textbf{Executive status.} P versus NP remains open.  TESTS-0020 correctly
computes a nonconstant four-interface semantic profile on two equal-size
signed-clean blocks and correctly constructs a total SAT wrapper that omits
that profile on their all-YES product island.  This run proves the exact
quantifier boundary: such island avoidance does not negate the worst-case
statement needed for a lower bound.  The negation of
$\forall A\,\exists F$ requires one $A$ avoiding the charged obligation on
\emph{all} inputs.  Avoidance only on a proper subfamily is compatible with
every decider having a charged witness outside it.  The profile calculation
is retained; any use of the easy island to retire the active worst-case
two-way transfer is rejected.}}
\end{center}

\section{Claim card}

\begin{center}\small
\begin{tabular}{@{}p{0.44\linewidth}p{0.20\linewidth}p{0.28\linewidth}@{}}
\toprule
Claim & Status & Scope\\ \midrule
\texttt{EASY-ISLAND-QUANTIFIER-GUARD1} & \EXACT & Arbitrary deciders, sources, islands, and obligations.\\
TESTS-0020 profile histograms & \EXACT & Finite $K_{3,3}$/prism blocks; review pending.\\
\texttt{EASY-ISLAND-REFUTES-WORSTCASE1} & \RETIRED & Invalid quantifier promotion.\\
Finite relation audit & \COMP & 1,043,584 relation/island pairs.\\
\texttt{QUANTITATIVE-TWO-WAY-ONLINE-TRANSFER1} & \OPEN & Same-machine worst-case transfer.\\
P versus NP & \OPEN & No unrestricted lower bound.\\
\bottomrule
\end{tabular}
\end{center}

\section{Exact quantifier theorem}

Let $X$ be the full encoded input language, let $I\subsetneq X$ be a
recognizable easy island, let $\mathcal D$ be a nonempty class of total exact
deciders, and let $R(A,x)$ be any proposed charged-obligation predicate.  No
computability or monotonicity of $R$ is needed for the logical theorem.

\begin{theorem}[\texttt{EASY-ISLAND-QUANTIFIER-GUARD1}]
Define
\begin{align}
 E_I &:\Longleftrightarrow
   \exists A\in\mathcal D\ \forall x\in I:\ \neg R(A,x), \tag{1}\\
 W_X &:\Longleftrightarrow
   \forall A\in\mathcal D\ \exists x\in X:\ R(A,x). \tag{2}
\end{align}
Then $E_I$ and $W_X$ are logically compatible.  In fact,
\[
 \neg W_X\quad\Longleftrightarrow\quad
 \exists A\in\mathcal D\ \forall x\in X:\ \neg R(A,x). \tag{3}
\]
If both (1) and (2) hold, every witness $A$ for (1) has some charged input in
$X\setminus I$.
\end{theorem}

\begin{proof}
Equation (3) is direct quantifier negation.  If $A$ witnesses (1) and (2)
holds, choose $x\in X$ with $R(A,x)$.  Such $x$ cannot lie in $I$ by (1), so
$x\in X\setminus I$.  Compatibility is witnessed, for example, by
$R(A,x)\Longleftrightarrow x\notin I$ for every $A$.
\end{proof}

\begin{corollary}[Application to TESTS-0020]
Let $I$ be the canonical disjoint products of the paired-incidence
$K_{3,3}$ and prism blocks.  The recognize-first wrapper $W_D$ proves that a
chosen semantic profile need not be computed on $I$.  It does not refute any
claim of form (2) on the full signed-clean family or on all encoded SAT
instances.  Therefore \texttt{EASY-ISLAND-REFUTES-WORSTCASE1} is \RETIRED.
\end{corollary}

This is not a criticism of the explicit scope statement in TESTS-0020: that
report already says its theorem ``does not refute an existential worst-case
lower bound.''  The red-team correction is that synthesis and future root
MAIN updates must preserve that restriction.  What is new in TESTS-0020 is
the exact histogram
\[
 \Pi_4(K_{3,3})=(36,9,81),\qquad
 \Pi_4(P)=(36,21,69) \tag{4}
\]
at projection cardinalities $(12,14,16)$.  The general wrapper conclusion is
an instance of earlier REDTAG1, representation-robustness, and source-
specificity warnings, not a new worst-case obstruction theorem.

\section{Reproducible hostile audit}

The dependency-free program
\begin{verbatim}
timeout 30s python3 experiments/review/easy_island_quantifier_guard_audit.py
\end{verbatim}
completed in 7.992623 seconds without randomness.  It exhausted every
obligation relation on three abstract deciders and two through five sources,
against every nonempty proper island.  It checked 1,043,584 relation/island
pairs.  In 405,604 cases easy-island avoidance and the worst-case obstruction
were simultaneously true.  All quantifier-negation and outside-witness checks
passed.

The named controls included the $K_{3,3}$ and prism clean blocks, Horn, 2SAT,
affine/Tseitin, bounded treewidth, independent products, padding, redundant
clauses, alternate encodings, and UNSAT.  DIRECT, COPY, DEAD, DELAY, INTERP,
and RECOMP wrappers all admitted the explicit compact predicate
$R(A,x)=[x\notin I]$.  A constant-size circuit or branching program computes
that predicate once island membership is provided; compact representation
does not change the quantifiers.  These finite checks are \COMP; Theorem 1 is
\EXACT.

Script SHA-256:
\begin{verbatim}
23d5d399f54f7a7c9d08a059e79b7d182fd371c7cd201d16d925207814aabc62
\end{verbatim}
Result SHA-256:
\begin{verbatim}
8fe7e1e7bb8d37dadcc3243ca0e7e1f059e303e68e271a8d082fa5008ac77d5a
\end{verbatim}

\section{Endpoint, barriers, and model matching}

FLMV proves a worst-case class exclusion: for each fixed $c$ below the golden
ratio, some $d>0$ forbids a deterministic random-access SAT machine that
meets both bounds on all inputs \cite{FLMV05}.  Consequently a compiler meant
to land in FLMV must establish the analogue of (2) in the same ordinary
machine realization.  Exhibiting a fast or profile-free proper subfamily does
not contradict that endpoint.

\paragraph{SAME\_REALIZATION.}
TESTS-0020's recognizer and fallback are one total machine, but its proof
accounts only for the island branch.  A worst-case transfer must also charge
the fallback path, every preprocessing pass, source retention, challenge,
query, replay, reconstruction, output, time, and peak workspace.  The unknown
polynomial exponent of the assumed SAT decider cannot be hidden in that
fallback.

\paragraph{Relativization, algebrization, and Natural Proofs.}
The quantifier theorem and its counterexample relativize and algebraize: the
same argument works with semantic or algebraic-oracle deciders.  It therefore
guards scope but supplies no escape from Baker--Gill--Solovay or
Aaronson--Wigderson \cite{BGS75,AW09}.  It defines no large constructive
circuit property, so Natural Proofs gives no positive inference.  Nonuniform
circuits do not repair the quantifier gap; (3) remains the required negation
for each declared decider class.

\paragraph{Post-quantum and cross-domain transfer.}
Schur powers, low-rank quadratic relations, syzygies, collision entropy, and
special-soundness forks may distinguish or bind a source.  That is an
\textbf{EXACT OBJECT/THEOREM TRANSFER} only for the declared representation.
Promoting a distinction on $I$ to (2) is a \textbf{FAILED TRANSFER}.
Worst-case/average-case architecture, observability, VLSI routing,
communication cuts, and proof-system search are \textbf{useful structural
analogies requiring a new transfer theorem}; their restricted resource is not
automatically the state of an arbitrary SAT decider.  Hidden algebraic
structure can make the island easier, as TESTS-0020 itself demonstrates.

\section{CORPUS REVIEWED}

The current GitHub corpus was surveyed using the previously enumerated 2,208
blob tree and fresh reads after TESTS-0020 publication.  Root
\path{main.tex}, \path{AGENTS.md}, \path{CODEX.md}, actual
\path{main4.tex}/\path{main5.tex}, the complete pnp-research skill/protocol,
and all task specifications were reviewed.  The newest verified
MAIN/SYNTHESIS-0020, predecessor TESTS-0020, LB-0023, TESTS-0018/0019,
REDTAG1, ORDERLOC1, ONETRANS1, PADQUOT1, FIBERDICH1, and their cited
experiments/results were read or reconciled by repository search.  Relevant
historical MAIN, governance, HEC/VSC-HEC/G-PPC-HEC/CRS-PPC/CCSPC, wrapper,
padding, source-specificity, response-fiber, and recognizable-island material
was inspected.  TESTS-0020 exists at its canonical path, is substantive, and
has the required header and final sentinel; no predecessor recovery request
was needed.

Primary literature was checked before the theorem: FLMV's original JACM
paper states the model-matched worst-case time--space exclusion
\cite{FLMV05}; Baker--Gill--Solovay and Aaronson--Wigderson provide the
relativization and algebrization barriers \cite{BGS75,AW09}; D\"ocker provides
the NP-complete signed clean family, not a claim that every subfamily is hard
\cite{Docker19}.

\section{DUPLICATION / CONTRADICTION AUDIT}

TESTS-0020 explicitly contains the correct caveat, so Theorem 1 confirms and
formalizes its scope rather than contradicting its finite construction.
REDTAG1 already used the same $K_{3,3}$/prism pair to show a satisfiable tag
can expose a promise answer while total-language robustness survives.
SYNTHESIS-0013 already showed answer-balanced cofinal layers do not force
pointwise extraction.  SYNTHESIS-0020 already showed source specificity is
not hardness.  The new contribution is the exact negation/compatibility
theorem (1)--(3) attached to the newest easy-island result; the profile values
in (4) remain TESTS-0020's new finite fact.

\section{WHAT CHANGED RELATIVE TO ROOT MAIN}

Root MAIN's active target is a worst-case quantitative two-way online
transfer.  That target is unchanged and remains \OPEN.  The LSPV-3 guard now
records that no proper easy island, even one with a nonconstant semantic
profile and a total recognize-first wrapper, can retire the target without a
single full-domain avoiding decider as in (3).

\section{ADOPTED / REJECTED / RETIRED PRIOR COMPONENTS}

Adopted: TESTS-0020's signed-clean membership proof, exact histograms, product
factorization, and same-machine easy branch; FLMV's exact worst-case machine
model; SYNTHESIS-0020's two-way target.  Rejected: interpreting a proper
island counterexample as the negation of a full-domain worst-case theorem.
Retired: \texttt{EASY-ISLAND-REFUTES-WORSTCASE1}.  Retained \OPEN:
\texttt{QUANTITATIVE-TWO-WAY-ONLINE-TRANSFER1}.

\section{STALE OR OVERSTATED CLAIMS FOUND}

Any future summary saying TESTS-0020 ``falsifies universal compatibility
extraction'' without specifying ``on every input of its easy island'' is
overstated.  The correct conclusion is that semantic-profile nonconstancy is
not sufficient for a pointwise obligation on that island.  It does not imply
the active worst-case transfer is false.  Counting semantic states or models
still gives no runtime bound.

\section{Precise handoff}

Preserve (3) as a mandatory quantifier check.  To falsify the active target,
construct one uniform polynomial SAT decider that avoids the claimed charged
object on the full input domain.  To prove the target, show that every assumed
polynomial SAT decider has at least one source outside all recognizable easy
islands where a fixed compiler forces both directions of a source-specific
online relation, with solver-exponent-independent simultaneous bounds and an
FLMV-forbidden SAME\_REALIZATION landing.  This remains \OPEN.

\begin{thebibliography}{9}
\bibitem{FLMV05}
L. Fortnow, R. Lipton, D. van Melkebeek, and A. Viglas,
``Time-space lower bounds for satisfiability,'' \emph{J. ACM} 52(6),
835--865, 2005. \url{https://doi.org/10.1145/1101821.1101822}.
\bibitem{BGS75}
T. Baker, J. Gill, and R. Solovay, ``Relativizations of the $P=?NP$
question,'' \emph{SIAM J. Comput.} 4(4), 431--442, 1975.
\url{https://doi.org/10.1137/0204037}.
\bibitem{AW09}
S. Aaronson and A. Wigderson, ``Algebrization: A new barrier in complexity
theory,'' \emph{ACM Trans. Comput. Theory} 1(1), 2009.
\url{https://eccc.weizmann.ac.il/report/2008/005/}.
\bibitem{Docker19}
J. D\"ocker, ``Monotone 3-Sat-$(2,2)$ is NP-complete,''
\emph{arXiv:1912.08032}, 2019.
\url{https://arxiv.org/abs/1912.08032}.
\end{thebibliography}

\noindent\textbf{Final status:} P versus NP remains \OPEN.

% END_OF_VERIFIED_REPORT::LSPV-3::QGUARD1
\end{document}
